\documentclass[10pt,a4paper]{amsart}
\usepackage[margin=19mm]{geometry}
\usepackage{amsmath,amssymb,mathtools}
\usepackage[scr=rsfs,cal=euler]{mathalfa}
\usepackage{xfrac}
\usepackage[unicode,hidelinks,bookmarks=false]{hyperref}
\newcommand{\ie}{i.e.\ }
\newcommand{\cf}{cf.\ }
\newcommand{\resp}{respectively\ }
\newcommand{\Subsection}{\subsection}
\providecommand{\nobreakdash}{\nobreak\hbox{-}}
\setlength{\emergencystretch}{2em}
\multlinegap0pt
\usepackage{amssymb}
\newtheorem{theorem}{Theorem}
\newtheorem{claim}{Claim}
\newcommand{\EE}{\mathbb{E}}
\newcommand{\PP}{\mathbb{P}}
\title{Determining a Point Configuration from a Subset of the Pairwise Distances}
\author{Itai Benjamini, Elad Tzalik}
\date{Source: arXiv:2208.13855v4 (CC BY 4.0)}
\begin{document}
\maketitle

\noindent\emph{Source and licence.} Determining a Point Configuration from a Subset of the Pairwise Distances. Version arXiv:2208.13855v4 (CC BY 4.0), distributed by arXiv under the Creative Commons Attribution 4.0 International licence, http://creativecommons.org/licenses/by/4.0/, as recorded in the arXiv licence metadata for this exact version and shown on the abstract page https://arxiv.org/abs/2208.13855v4. Full licence notice: \texttt{licenses/arxiv-2208.13855-CC-BY-4.0.txt}.\par\smallskip

\noindent\textit{Editorial scope.} Counted unit: the article's theorem thm:monotone-threshold (source0.tex lines 866--872). The article's complete original proof of it is delivered with it (source0.tex lines 874--1014, one \texttt{proof} environment of 141 lines). No mathematics is written, repaired or re-derived: the statement and the proof are the article's own lines, in the article's own notation, and no retained line is substituted or re-flowed.  No further source-local context is needed: every cross-reference of every retained line resolves inside the two delivered spans.  Nothing was omitted from any retained span, so the package carries no omission record.  No retained line contains a citation, so the wrapper carries no bibliography and no bibliography entry.  No retained line contains a figure environment, an image-inclusion command or drawing code, so no image is delivered and the package declares no asset dependency.  Editorial formatting only, in addition: an AMS \texttt{amsart} preamble that restates the article's own declarations and macros used by the retained lines, namely the environments \texttt{theorem}, \texttt{claim} on separate counters, together with the standard packages those lines need.  The retained environments print in this document as Theorem 1, Claim 1 in order of appearance, whereas the article prints the same items with the numbering of its own full document.  The complete native source of the article as distributed in the arXiv e-print for this exact version is bundled unchanged as \texttt{sources/130051/source0.tex}.
\begin{theorem}\label{thm:monotone-threshold}
Let \(G \sim G(n,p)\) be a random graph on \([n]\), with its natural order. Then:
\begin{itemize}
    \item for every \(0<\varepsilon<1\), if \(p=(1-\varepsilon)\frac{\ln n}{n}\), then w.h.p. there is no monotone path from \(1\) to \(n\);
    \item for every \(\varepsilon>0\), if \(p=(1+\varepsilon)\frac{\ln n}{n}\), then w.h.p. there is a monotone path from \(1\) to \(n\).
\end{itemize}
\end{theorem}

\begin{proof}
Let \(N\) be the number of monotone paths from \(1\) to \(n\). A monotone path
with \(k\) internal vertices is determined by choosing \(k\) vertices from
\(\{2,\ldots,n-1\}\), so
\(\EE N=\sum_{k=0}^{n-2}\binom{n-2}{k}p^{k+1}=p(1+p)^{n-2}\le p e^{pn}\).
For \(p=(1-\varepsilon)\ln n/n\), this is \(o(1)\), and Markov's inequality
concludes the subcritical case.

\noindent It remains to analyze the supercritical case.
Write \(p=c\ln n/n\), where \(c=1+\varepsilon>1\).

\begin{claim}\label{clm:one-sided-growth}
Let \(p=c\ln n/n\), where \(c>1\). Then there is \(\tau>1/2\), depending
on \(c\), such that w.h.p. at least \(n^\tau\) vertices in
\(\{1,\ldots,\lfloor n/2\rfloor\}\) are reachable from \(1\) by monotone paths
contained in this set.
\end{claim}

\begin{proof}[Proof of the claim]
Choose \(1<c_0<c\), then choose \(K\) so large that
\(r:=K\ln(1+c_0/K)>1\); this is possible since \(r\to c_0\). Choose
\(\delta>0\) so small that
\[
    \left(\frac12-2\delta\right)r>\frac12,
\]
and then fix
\[
    \frac12<\tau<\min\left\{1,\left(\frac12-2\delta\right)r\right\}.
\]

Let \(I_0=\{2,\ldots,\lfloor\delta n\rfloor\}\). Expose the edges from \(1\) to
\(I_0\). Since \(\EE |N(1)\cap I_0|=(c\delta+o(1))\ln n\), Chernoff gives,
for \(a=c\delta/2\), w.h.p.
\[
    |N(1)\cap I_0|\ge a\ln n .
\]
On this event, put \(S_0=N(1)\cap I_0\), and observe that all vertices in
\(S_0\) are connected to \(1\) by a monotone edge. Set
\[
    M=\left\lfloor \frac{n}{K\ln n}\right\rfloor,
    \qquad
    T=\left\lfloor\left(\frac12-2\delta\right)K\ln n\right\rfloor .
\]
Let \(J_1,\ldots,J_T\) be consecutive intervals of length \(M\) immediately
after \(I_0\). As \(n\) is large, they all lie in
\(\{1,\ldots,\lfloor n/2\rfloor\}\), since
\[
    \lfloor\delta n\rfloor+TM
    \le \delta n+\left(\frac12-2\delta\right)n
    <\frac n2 .
\]

Let \(m=\lceil n^\tau\rceil\). We process the \(J_t\)'s in order, maintaining
sets \(S_t\) of vertices reachable from \(1\) by monotone paths contained in
\(I_0\cup J_1\cup\cdots\cup J_t\). Once \(|S_{t-1}|\ge m\), set
\(S_t=S_{t-1}\). Otherwise expose the edges between \(S_{t-1}\) and \(J_t\),
and set
\[
    S_t=S_{t-1}\cup\{y\in J_t: N(y)\cap S_{t-1}\ne\emptyset\}.
\]
Every vertex of \(S_t\) is reachable from \(1\) by a monotone path, since
\(J_t\) lies to the right of all previously used intervals.

Let \(\mathcal F_{t-1}\) be the history before step \(t\), and suppose
\(|S_{t-1}|=x<m\). Then \(x\ge a\ln n\). No edge from \(S_{t-1}\) to \(J_t\)
has previously been exposed, so conditional on \(\mathcal F_{t-1}\),
\[
    |S_t\setminus S_{t-1}|
    \sim \operatorname{Bin}\left(M,1-(1-p)^x\right).
\]
As \(\tau<1\), uniformly for \(x<m\) we have \(px=o(1)\), and therefore
\[
    \EE\bigl(|S_t\setminus S_{t-1}|\mid \mathcal F_{t-1}\bigr)
    =\left(\frac{c}{K}+o(1)\right)x .
\]
Since \(c_0<c\), Chernoff gives, for some constant \(\gamma>0\),
\[
    \PP\left(|S_t\setminus S_{t-1}|<\frac{c_0}{K}x
    \mid \mathcal F_{t-1}\right)
    \le e^{-\gamma x}
    \le n^{-\gamma a}.
\]
Call a step \(t\) active if \(|S_{t-1}|<m\), and bad if it is active and
\[
    |S_t\setminus S_{t-1}|<\frac{c_0}{K}|S_{t-1}|.
\]
For each \(t\), the conditional probability that \(t\) is bad, given
\(\mathcal F_{t-1}\), is at most \(n^{-\gamma a}\). Hence
\[
    \PP(\text{some bad step})\le Tn^{-\gamma a}=o(1).
\]
On the complementary event, every active step satisfies
\[
    |S_t|\ge \left(1+\frac{c_0}{K}\right)|S_{t-1}|.
\]
Thus, if the process has not reached \(m\) by time \(T\), then all \(T\) steps
were active and none was bad, so
\[
    |S_T|
    \ge a\ln n\left(1+\frac{c_0}{K}\right)^T
    =n^{(1/2-2\delta)r+o(1)}
    > n^\tau
\]
for all large \(n\). Since \(|S_T|\) is an integer, this implies
\(|S_T|\ge \lceil n^\tau\rceil=m\), a contradiction. Therefore, with
probability \(1-o(1)\), the process reaches size \(m\), and in particular
\(|S_T|\ge n^\tau\), proving the claim.
\end{proof}

\noindent To complete the proof, apply Claim~\ref{clm:one-sided-growth} in the left
half. W.h.p. there is a set
\[
    L\subseteq \{1,\ldots,\lfloor n/2\rfloor\},
    \qquad |L|\ge n^\tau,
\]
such that every \(u\in L\) is reachable from \(1\) by a monotone path contained
in the left half.

Applying the same argument to the right half in the reverse order, w.h.p.
there is a set
\[
    R\subseteq \{\lfloor n/2\rfloor+1,\ldots,n\},
    \qquad |R|\ge n^\tau,
\]
such that every \(v\in R\) has a monotone path from \(v\) to \(n\) contained in
the right half.

The two explorations expose only edges inside the two halves, so all edges
between \(L\) and \(R\) are still unexposed. Conditional on \(L\) and \(R\),
\[
    \PP(e(L,R)=0)
    \le (1-p)^{|L||R|}
    \le \exp(-p n^{2\tau})
    =o(1),
\]
because \(\tau>1/2\). Hence w.h.p. there is an edge \(uv\) with
\(u\in L\), \(v\in R\). Since \(u<v\), concatenating a monotone path from
\(1\) to \(u\), the edge \(uv\), and a monotone path from \(v\) to \(n\), gives
a monotone path from \(1\) to \(n\). This proves the supercritical assertion,
and hence the theorem.
\end{proof}

\end{document}
